\documentclass[11pt,a4paper]{article}
\usepackage[utf8]{vietnam}
\usepackage{amsmath,amssymb}
\usepackage{tabularx}
\usepackage{tikz}
\usepackage{xcolor}
\usepackage{fancyhdr}
\usepackage[top=1.8cm, bottom=2.0cm, left=1.8cm, right=1.8cm]{geometry}

\setlength{\parskip}{3pt}
\setlength{\parindent}{0pt}

\pagestyle{fancy}
\fancyhf{}
\renewcommand{\headrulewidth}{0.4pt}
\lhead{\footnotesize \textbf{CHÙA HƯƠNG SƠN -- HƯỚNG DẪN CHẤM THI}}
\rhead{\footnotesize \textbf{ĐÁP ÁN CHI TIẾT ĐỀ THI THỬ SỐ 04}}
\cfoot{\thepage}

\begin{document}

\begin{center}
{\large\bfseries ĐÁP ÁN \& HƯỚNG DẪN CHẤM ĐỀ THI THỬ CUỐI HỌC KỲ I (ĐỀ SỐ 04)}\\[1.5mm]
{\bfseries ĐƠN VỊ: CHÙA HƯƠNG SƠN \quad|\quad MÔN: TOÁN -- LỚP 9 \quad|\quad NĂM HỌC 2026 -- 2027}\\[1mm]
\textit{(Chuẩn cấu trúc định dạng mới của Sở Giáo dục và Đào tạo)}
\end{center}

\vspace{1mm}
\hrule
\vspace{3mm}

\noindent
\textbf{\large PHẦN I. TRẮC NGHIỆM KHÁCH QUAN (7,0 ĐIỂM)}

\vspace{2mm}
\noindent
\textbf{1. Bảng đáp án Phần 1: Trắc nghiệm nhiều lựa chọn (3,0 điểm)} \textit{(Mỗi câu đúng được 0,25 điểm)}
\begin{center}
\begin{tabular}{|c|c|c|c|c|c|c|c|c|c|c|c|c|}
\hline
\textbf{Câu} & \textbf{1} & \textbf{2} & \textbf{3} & \textbf{4} & \textbf{5} & \textbf{6} & \textbf{7} & \textbf{8} & \textbf{9} & \textbf{10} & \textbf{11} & \textbf{12} \\
\hline
\textbf{Đáp án} & \textbf{A} & \textbf{A} & \textbf{A} & \textbf{A} & \textbf{A} & \textbf{A} & \textbf{A} & \textbf{A} & \textbf{A} & \textbf{A} & \textbf{A} & \textbf{A} \\
\hline
\end{tabular}
\end{center}

\vspace{2mm}
\noindent
\textbf{2. Bảng đáp án Phần 2: Trắc nghiệm Đúng -- Sai (2,0 điểm)}
\begin{center}
\begin{tabular}{|c|c|c|c|c|l|}
\hline
\textbf{Câu} & \textbf{Ý a} & \textbf{Ý b} & \textbf{Ý c} & \textbf{Ý d} & \textbf{Ghi chú đáp án chuẩn} \\
\hline
\textbf{Câu 13} & \textbf{Đ} & \textbf{Đ} & \textbf{S} & \textbf{Đ} & Ý c) sai vì vế phải bằng $0{,}25 \times 400 = 100\text{ gam}$. \\
\hline
\textbf{Câu 14} & \textbf{Đ} & \textbf{Đ} & \textbf{Đ} & \textbf{S} & Ý d) sai vì tích $AC \cdot BD = R^2$ (chứ không phải $2R^2$). \\
\hline
\end{tabular}
\end{center}

\vspace{2mm}
\noindent
\textbf{3. Bảng đáp án Phần 3: Trắc nghiệm trả lời ngắn (2,0 điểm)} \textit{(Mỗi câu đúng được 0,5 điểm)}
\begin{center}
\begin{tabular}{|c|l|c|c|}
\hline
\textbf{Câu} & \textbf{Nội dung câu hỏi} & \textbf{Kết quả bài làm} & \textbf{Điểm} \\
\hline
\textbf{Câu 15} & Nghiệm bất phương trình $5(x + 1) - 3 > 3(2x - 1) + 4$ & \textbf{$x < 1$} & 0,5đ \\
\hline
\textbf{Câu 16} & Giá trị biểu thức $K = \sqrt{(\sqrt{8}-3)^2} + \sqrt{8}$ & \textbf{$3$} & 0,5đ \\
\hline
\textbf{Câu 17} & Độ cao đỉnh thang chạm tường ($4 \cdot \sin 65^\circ$) & \textbf{$3{,}63\text{ m}$} & 0,5đ \\
\hline
\textbf{Câu 18} & Diện tích miếng kim loại ($\dfrac{3{,}14 \times 25^2 \times 72}{360}$) & \textbf{$392{,}5\text{ cm}^2$} & 0,5đ \\
\hline
\end{tabular}
\end{center}

\vspace{3mm}
\noindent
\textbf{\large PHẦN II. TỰ LUẬN (3,0 ĐIỂM)}

\vspace{2mm}
\textbf{Câu 19 (1,0 điểm).}
\begin{itemize}
  \item[a)] \textit{(0,5 điểm)} $\begin{cases} 5x + 2y = 9 \\ 2x - y = 0 \end{cases} \Leftrightarrow \begin{cases} 5x + 2(2x) = 9 \\ y = 2x \end{cases} \Leftrightarrow \begin{cases} 9x = 9 \\ y = 2x \end{cases} \Leftrightarrow \begin{cases} x = 1 \\ y = 2 \end{cases}$. \hfill \textbf{(0,5đ)}
  \item[b)] \textit{(0,5 điểm)} $M = \dfrac{1}{\sqrt{x}(\sqrt{x}-1)} : \dfrac{(\sqrt{x}+1)(\sqrt{x}-1) - (\sqrt{x}+2)(\sqrt{x}-2)}{(\sqrt{x}-2)(\sqrt{x}-1)} = \dfrac{1}{\sqrt{x}(\sqrt{x}-1)} : \dfrac{3}{(\sqrt{x}-2)(\sqrt{x}-1)} = \dfrac{\sqrt{x}-2}{3\sqrt{x}}$. \hfill \textbf{(0,5đ)}
\end{itemize}

\vspace{1mm}
\textbf{Câu 20 (1,5 điểm).}
\begin{itemize}
  \item[a)] \textit{(0,75 điểm)} Theo tính chất hai tiếp tuyến cắt nhau: $OC$ là phân giác $\widehat{AOM}$, $OD$ là phân giác $\widehat{BOM} \Rightarrow \widehat{COD} = \dfrac{1}{2}(\widehat{AOM} + \widehat{BOM}) = 90^\circ$. Lại có $CA = CM, DB = DM \Rightarrow CD = CM + MD = AC + BD$. \hfill \textbf{(0,75đ)}
  \item[b)] \textit{(0,75 điểm)} Ta có $AC \parallel BD \Rightarrow \dfrac{AK}{KD} = \dfrac{AC}{BD} = \dfrac{CM}{MD} \Rightarrow MK \parallel AC$. Theo định lí Thales: $\dfrac{MK}{AC} = \dfrac{MD}{CD} \Rightarrow MK = \dfrac{AC \cdot MD}{CD} = \dfrac{AC \cdot BD}{AC + BD} \Rightarrow MK \cdot (AC + BD) = AC \cdot BD$. \hfill \textbf{(0,75đ)}
\end{itemize}

\vspace{1mm}
\textbf{Câu 21 (0,5 điểm).}
\begin{itemize}
  \item[a)] \textit{(0,25 điểm)} Cạnh đáy hộp là $60 - 2x$ (cm, $0 < x < 30$). Thể tích $V = x(60 - 2x)^2 = \dfrac{1}{4} \cdot 4x(60 - 2x)(60 - 2x)$. Áp dụng BĐT AM-GM cho ba số dương:
  $4x(60 - 2x)(60 - 2x) \le \left(\dfrac{4x + 60 - 2x + 60 - 2x}{3}\right)^3 = 40^3 = 64000 \Rightarrow V \le 16000\text{ cm}^3$. Dấu ``='' khi $4x = 60 - 2x \Leftrightarrow x = 10\text{ cm}$. \hfill \textbf{(0,25đ)}
  \item[b)] \textit{(0,25 điểm)} Ta có $S = \dfrac{1}{x^2 + y^2} + \dfrac{1}{2xy} + \dfrac{1}{2xy} \ge \dfrac{4}{(x+y)^2} + \dfrac{1}{2 \cdot \frac{(x+y)^2}{4}} = \dfrac{4}{(x+y)^2} + \dfrac{2}{(x+y)^2} = \dfrac{6}{(x+y)^2} \ge \dfrac{6}{2^2} = \dfrac{3}{2}$. Dấu ``='' khi $x = y = 1$. Vậy $\min S = \dfrac{3}{2}$. \hfill \textbf{(0,25đ)}
\end{itemize}

\end{document}
